\section{Density centers survive the height blow-up}\label{sec:density-centers}

The measures of Section~\ref{sec:linearized} retain height and tilt, but
the final dimension argument also needs the large set of density-$Q$
centers. Ordinary monotonicity about those centers is not accurate enough
at the height-square scale: a curvature error depending only on the
reference graph could dominate $a_i^2$. We remove that error by measuring
radius along the reference graph and compensating its volume density.

Fix the integer $p\ge4$ chosen above. For $x,y\in\Pi$ and $l>0$, set
\begin{equation}\label{eq:limit-kernels}
 \phi_{x,l}(y)=l^{-m}\left(1-\frac{|y-x|^2}{l^2}\right)_+^p,
 \qquad
 \psi_{x,l}(y)=2p\,l^{-m}
                  \left(1-\frac{|y-x|^2}{l^2}\right)_+^{p-1}.
\end{equation}

\begin{proposition}[Comparison at the density centers]
\label{prop:density-comparison}
In the setting of Proposition~\ref{prop:linearized}, suppose that
$\Omega$ contains a fixed neighborhood of $\overline B_1\subset\Pi$.
Let $\sigma_i$ be positive measures supported on density-$Q$ points of
$V_i$ in the ambient closed unit ball, with
\[
 c\le\sigma_i(\overline{\mathbf B}_1)\le C,
 \qquad \sigma_i(\mathbf B_r(X))\le Cr^s
 \quad(X\in\R^{m+n},\ r>0),\qquad s>m-1.
\]
Pass to a weak limit $\sigma$ along the subsequence defining the
linearized measures. Then $\sigma\ne0$, its support lies in
$\overline B_1\subset\Pi$, and it has the same upper ball bound.
There is $l_*>0$ such that
\begin{equation}\label{eq:center-comparison}
 2l^2\int\phi_{x,l}\,\dd\mathfrak m
       \ge\int\psi_{x,l}\,\dd\nu
 \qquad(x\in\spt\sigma, 0<l<l_*).
\end{equation}
\end{proposition}

\begin{proof}
Weak compactness and the stated conclusions about $\sigma$ follow from
the mass and ball bounds and from convergence of the varifold supports
to $\Pi$. We prove \eqref{eq:center-comparison} for each fixed
$0<l<l_*$, where $l_*$ is chosen so all the following constructions
have common interior room. All estimates below are uniform over
$X_0\in\spt\sigma_i$.

\subsection*{A radial test adapted to the reference graph}

Write $w=X_0-\pi_iX_0$ and translate the reference graph to
$S=S_i+w$, which passes through $X_0$. By the height bound,
$|w|\le Ca_i$. In this part of the proof, $\pi,z,B,e,J_z$ refer to this
translated graph, and $\rho$ is geodesic distance on $S$ from $X_0$,
pulled back by $\pi$. Translation preserves its geometric estimates.
Changing the reference tangent plane by this translation and the nearby
projection changes squared tilt by at most
\begin{equation}\label{eq:translated-tilt}
                  e\le C(e_i+a_i^2).
\end{equation}
Indeed the two projections differ by $O(a_i)$ and the second fundamental
forms are uniformly bounded; squared tilt is comparable to the squared
distance of the corresponding orthogonal projections.

Define
\[
 d^2=\rho^2+|z|^2,\qquad
 U=\rho\nabla_S\rho,\qquad
 G=\nabla(d^2/2)=J_z^{-1}U+z,\qquad Y=J_zU+z.
\]
These fields are smooth at the central fiber, and
$Y\cdot G=d^2$. Let $W$ be the inverse volume density in geodesic
normal coordinates on $S$ centered at $X_0$, again pulled back by
projection. Thus integration of $W\vol_S$ is ordinary integration in
those coordinates. Enlarging the sequence $\varepsilon_i\to0$ from
\eqref{eq:horizontal-error-bound} if necessary, geodesic coordinate
estimates give
\begin{equation}\label{eq:geodesic-compensation}
\begin{gathered}
 |\nabla_S W|\le C\varepsilon_i\rho,
 \qquad |\nabla^S U-\Id|\le C\varepsilon_i\rho^2,\\
 W(\diver_SU-m)+U\cdot\nabla_SW=0.
\end{gathered}
\end{equation}
The last identity is exact: $U$ is the coordinate radial field and $W$
is the reciprocal volume density.

Apply \eqref{eq:horizontal-error} to $U$ and add the normal trace
\eqref{eq:normal-trace}. Their leading tilt terms cancel. The identity
in \eqref{eq:geodesic-compensation} then cancels the reference-volume
error. The result is
\begin{equation}\label{eq:radial-error}
\begin{aligned}
 \mathcal E&=P_T:D(WY)-mW,\\
 |\mathcal E|&\le C\varepsilon_i\bigl\{
 e(\rho^2+|z|)+(a_i|z|+|z|^2)(1+\rho)
                         +\rho\sqrt e\,|z|\bigr\}.
\end{aligned}
\end{equation}
For completeness, pairing $\nabla^SU-\Id$ with $B$ costs
$C\varepsilon_i\rho^2e$. Differentiating $W$ in the horizontal part
gives $U\cdot\nabla_SW$ and another error of this size. Its vertical
pairing costs $C\varepsilon_i\rho\sqrt e\,|z|$. The remaining terms
are exactly those bounded in \eqref{eq:horizontal-error-bound} and
\eqref{eq:normal-trace}.

Put
\[
 \Phi_t(d)=\left(1-\frac{d^2}{t^2}\right)_+^p,
 \qquad \Psi_t(d)=2p\left(1-\frac{d^2}{t^2}\right)_+^{p-1},
 \qquad
 I_i(X_0,t)=t^{-m}\int W\Phi_t(d)\,\dd\mu_i.
\]
Differentiate in $\log t$ and test stationarity with
$W\Phi_t(d)Y$. Since $Y\cdot G=d^2$, this gives
\begin{equation}\label{eq:radial-derivative}
 \dot I_i(X_0,t)
 =t^{-m}\int\left\{
       \frac{\Psi_t}{t^2}W Y\cdot(I-P_T)G
                         +\Phi_t\mathcal E\right\}\,\dd\mu_i.
\end{equation}
The test is compactly supported within the normal tube; the cutoff power
$p\ge4$ gives more regularity than this calculation requires.

Here the estimate must hold down to arbitrarily small $t$, with no
comparison between $t$ and $a_i$. On $d\le t\le l_*\le1$,
\begin{equation}\label{eq:radial-normal-control}
\begin{aligned}
 Y\cdot(I-P_T)G&\ge\tfrac12|(I-P_T)G|^2
                              -C\varepsilon_i^2|z|^2t^2,\\
 |z|&\le |(I-P_T)G|+Ct\sqrt e.
\end{aligned}
\end{equation}
The first inequality follows from
$|Y-G|\le C\varepsilon_i|z|\rho$. For the second, take the vertical
component of $G$ and use $|P_vP_T|\le C\sqrt e$ and $|G|\le Ct$.
Set $b'=|(I-P_T)G|/t$. Then both $|z|\le t$ and
$|z|\le t(b'+C\sqrt e)$ hold. In particular,
\[
 e|z|\le te,\qquad
 a_i|z|\le Ct\bigl((b')^2+e+a_i^2\bigr),
 \qquad |z|^2\le Ct^2\bigl((b')^2+e\bigr).
\]
Using $\Phi_t\le C\Psi_t$, every adverse term in
\eqref{eq:radial-derivative} is therefore at most
$C\varepsilon_i t\Psi_t((b')^2+e+a_i^2)$.
Absorb the $(b')^2$ term into the positive term from
\eqref{eq:radial-normal-control}. Equation~\eqref{eq:translated-tilt}
and comparison of $d$ with ambient distance yield
\begin{equation}\label{eq:radial-almost-monotonicity}
 \dot I_i(X_0,t)\ge
 -C\varepsilon_i t^{1-m}
       \int_{\mathbf B_{Ct}(X_0)}(e_i+a_i^2)\,\dd\mu_i.
\end{equation}

\subsection*{Averaging over the centers}

Let $c_\Phi=\int_{\R^m}(1-|y|^2)_+^p\,\dd y$.
At the center $X_0$, the squared radius $d^2$ agrees with the squared
ambient distance to leading order and $W(X_0)=1$. Since $X_0$ has
density $Q$,
\[
                  \lim_{t\downarrow0}I_i(X_0,t)=Qc_\Phi.
\]
Take any bounded nonnegative continuous center weight $g$. Average
\eqref{eq:radial-almost-monotonicity} against $g\,\dd\sigma_i$ and
integrate in logarithmic radius. Fubini's theorem and the upper ball
bound for $\sigma_i$ give
\begin{align}
 &\int g(X_0)\bigl(I_i(X_0,l)-Qc_\Phi\bigr)\,\dd\sigma_i(X_0)
 \notag\\
 &\hspace{2em}\ge
 -C\|g\|_\infty\varepsilon_i
    \left(\int(e_i+a_i^2)\,\dd\mu_i\right)
                  \int_0^l t^{s-m}\,\dd t
 =-o(a_i^2).\label{eq:averaged-density}
\end{align}
The first integral is over fixed interior room and is $O(a_i^2)$.
The second is finite precisely because $s>m-1$. To justify the lower
endpoint, first integrate between positive radii, then use Fatou's
lemma for the nonnegative functions $gI_i$ and the displayed density
limit. This avoids any uniform rate of density convergence in $X_0$.

The averaged inequality now gives a lower bound with error $o(a_i^2)$.
To obtain the asserted comparison, we need an upper bound for the same
integrand in terms of the original mass defect and squared height.
Translating the reference graph introduces a cross term and an additional
negative square in this upper expansion. The vanishing first moment
removes the cross term, and the negative square can be discarded.

\subsection*{Expansion at a fixed radius}

Keep $l>0$ fixed. Write $y=\pi_iX$ and $z_i=X-y$. After translating
the new base back to $S_i$, its projection and residual are
\[
                   y'=\pi_i(X-w),\qquad z'=X-w-y'.
\]
All relevant support points satisfy $|z_i|+|z'|\le Ca_i$ in common
interior room. Functions such as
$q(y)=W(y+w)\Phi_l(\rho(y+w))$, and the analogous function with
$\Psi_l$, have uniform $C^2$ bounds and compact interior support.
We claim that, uniformly in the centers,
\begin{equation}\label{eq:translated-base-test}
             \int q(y')\,\dd\mu_i
               =\int q(y)\,\dd\mu_i+o(a_i^2).
\end{equation}
Indeed the nearest-projection formula and smooth convergence to a plane
give
\[
 y'-y=-P_h(y)w+o(a_i^2),\qquad
                  \sup_y|P_h(y)w|=o(a_i).
\]
For the first estimate, $D\pi_i(X)=J_{z_i}^{-1}P_h(y)$ differs from
$P_h(y)$ by $O(\varepsilon_i a_i)$, and the second derivative of
$\pi_i$ is $o(1)$. For the second, $w$ is normal to $S_i$ at the
center and all tangent planes tend uniformly to $\Pi$.
Taylor expansion thus leaves only
$-\int\nabla^{S_i}q\cdot P_hw\,\dd\mu_i$, up to $o(a_i^2)$.
Replacing the projected measure by $Q\vol_{S_i}$ costs $o(a_i^2)$
because $|M_i|=O(a_i^2)$. Integration by parts then uses
\[
             \diver_{S_i}(P_hw)=\mathcal H_{S_i}\cdot w=o(a_i^2)
\]
to prove \eqref{eq:translated-base-test}. Compact support eliminates
all boundary terms.

Taylor expansion in the squared normal height gives, uniformly on the
test region,
\[
 \Phi_l\bigl((\rho^2+|z'|^2)^{1/2}\bigr)
 =\Phi_l(\rho)-\frac{|z'|^2}{2l^2}\Psi_l(\rho)+o(a_i^2).
\]
The normalized reference-volume integral of $W\Phi_l(\rho)$ is
exactly $c_\Phi$, by its geodesic-coordinate definition. Hence
\eqref{eq:translated-base-test} expresses the zeroth-order term of
$I_i-Qc_\Phi$ as an integral against $M_i$, with error $o(a_i^2)$.

For the height term, $z'=z_i-w+o(a_i)$ uniformly. Put
$q_\Psi(y)=W(y+w)\Psi_l(\rho(y+w))\ge0$. Moving its argument
from $y'$ to $y$ costs $o(a_i^2)$ in the height-square integral,
and
\[
 \int q_\Psi|z_i-w|^2\,\dd\mu_i
 =\int q_\Psi|z_i|^2\,\dd\mu_i
   +|w|^2\int q_\Psi\,\dd\mu_i
   -2w\cdot\int q_\Psi z_i\,\dd\mu_i.
\]
This integral enters $I_i$ with the negative coefficient
$-l^{-m}/(2l^2)$, so discarding the $|w|^2$ term gives an upper
bound. The mixed term is $o(a_i^2)$ by the centering assumption.
We spell out its uniformity, since the centers move with $i$.
The vector measures
$a_i^{-1}(\pi_i)_\#(z_i\mu_i)$ have uniformly bounded total variation
and converge weakly to zero. At fixed $l$, the above center-dependent
tests and their smooth coefficient factors form a relatively compact
family in the uniform norm on common interior room. A finite
approximation of this family therefore makes the weak convergence
uniform over all its members. Multiplication by the bounded vectors
$w/a_i$ preserves this conclusion; these vectors need not converge.

It follows that, after division by $a_i^2$, the limiting upper bound
for $I_i(X_0,l)-Qc_\Phi$ at a center tending to $x\in\Pi$ is
\begin{equation}\label{eq:center-expression}
 F_l(x)=\int\phi_{x,l}\,\dd\mathfrak m
                  -\frac1{2l^2}\int\psi_{x,l}\,\dd\nu.
\end{equation}
More explicitly, let $\operatorname{pr}_{\Pi}$ denote orthogonal
projection onto the limiting plane. The same compact-test argument
applied to the bounded measure sequences in \eqref{eq:limit-measures}
gives
\begin{equation}\label{eq:uniform-center-upper}
 a_i^{-2}\bigl(I_i(X_0,l)-Qc_\Phi\bigr)
 \le F_l(\operatorname{pr}_{\Pi}X_0)+o(1)
 \quad\hbox{uniformly for }X_0\in\spt\sigma_i.
\end{equation}
The function $F_l$ is continuous.
Combining this uniform upper bound with
\eqref{eq:averaged-density}, and passing to the limit against any
nonnegative continuous $g$, gives $\int gF_l\,\dd\sigma\ge0$.
Continuity implies $F_l\ge0$ on $\spt\sigma$, which is
\eqref{eq:center-comparison}. Since $l\in(0,l_*)$ was arbitrary, the
comparison holds for every such radius.
\end{proof}
